\section{Configuration-Robust Correctness}
\label{sec:metric}

We formalize the metrics used throughout the paper. The guiding idea is to
separate correctness that a developer (or an agent) can \emph{observe cheaply}---by
running a program once at low parallelism---from correctness that actually
\emph{holds} once the program executes across the parallel configurations a
production run will use.

\subsection{Setup and Notation}
Let $\mathcal{M}$ denote the set of models under evaluation and $\mathcal{T}$ a
set of parallel-programming tasks. Each task $t\in\mathcal{T}$ is specified by a
natural-language prompt with an explicit function contract and is equipped with a
\emph{trusted serial reference} $R_t$: a sequential implementation that defines
the correct output on any input. A model $m$ produces a program
\begin{equation}\label{eq:gen}
G(m,t) \;=\; \textsc{extract}\big(m(\textsc{prompt}(t))\big),
\end{equation}
where $\textsc{extract}(\cdot)$ recovers the code from the response. The generated
program and the reference are compiled into the same driver, which supplies
randomized inputs and adjudicates the program's output against $R_t$.

A parallel program admits a family of \emph{execution configurations}
$\mathcal{K}$. A configuration fixes a degree of parallelism---here an OpenMP
thread count $\tau\in\{1,2,4,8\}$---and a repetition index that re-runs the same
program to expose nondeterminism. Executing program $G$ under configuration
$k\in\mathcal{K}$ yields either a validated output or the failure symbol $\bot$
(a crash or a compilation failure), and the driver reports a per-configuration
correctness indicator
\begin{equation}\label{eq:corr}
\mathrm{corr}(G,k) \;=\;
\begin{cases}
1 & G\text{ runs under }k\text{ and its output matches }R_t\\
  & \text{on every validation input,}\\
0 & \text{otherwise.}
\end{cases}
\end{equation}
Correctness is checked against the serial reference on multiple randomized
inputs, so a program must be right on all of them, not merely on a fixed example.

\subsection{Single-Configuration vs.\ Robust Correctness}
Let $k_0\in\mathcal{K}$ denote the \emph{developer-default} configuration: a
single thread ($\tau{=}1$), the most forgiving setting and the one a quick local
check most naturally uses. Because a single-threaded run executes the parallel
constructs sequentially, it hides precisely the hazards---races,
non-associative combines, order dependence---that only manifest when the work is
actually divided. We define two correctness notions for a generated program $G$:
\begin{align}
\text{single-configuration:}\quad \mathcal{S}(G) &= \mathrm{corr}\big(G,k_0\big),\label{eq:S}\\
\text{configuration-robust:}\quad \mathcal{R}(G) &= \!\!\prod_{k\in\mathcal{K}}\!\! \mathrm{corr}(G,k).\label{eq:R}
\end{align}
$\mathcal{S}$ is what an ordinary single-run check observes; $\mathcal{R}$
requires agreement with the reference under \emph{every} thread count and every
repetition. We also define \emph{configuration-invariance} $\mathcal{I}(G)=1$ iff
$G$ produces identical outputs for all $k$ (regardless of whether they are
correct)---a program can be invariant yet uniformly wrong.

\subsection{Correctness Is Necessary, Not Sufficient: The Second Axis}
\label{sec:secondaxis}
$\mathcal{R}$ has a deliberate blind spot that is central to this paper: it is a
predicate on \emph{outputs}, so it is \emph{maximized by removing parallelism
altogether}. A program that deletes every \texttt{\#pragma}, or wraps the loop body
in one \texttt{critical}, agrees with the serial reference at every thread
count and is therefore configuration-robustly correct---while running no faster,
or slower, than serial code. A correctness metric for parallel code that a serial
program aces is not, by itself, measuring parallel code. We therefore pair
$\mathcal{R}$ with a \emph{performance} axis. For a program that compiles and
validates, let $T(G,p)$ be its wall-time on $p$ threads; its \emph{self-speedup}
is $\mathcal{P}_p(G)=T(G,1)/T(G,p)$ and its parallel efficiency
$\mathcal{P}_p(G)/p$. An acceptance test for agent-generated parallel code needs
\emph{both} axes: $\mathcal{R}(G){=}1$ (robustly correct) \emph{and}
$\mathcal{P}_p(G)$ meaningfully above $1$ (actually parallel). Section~\ref{sec:results}
shows the second axis, not the first, is where accepted code fails.

\subsection{$\mathcal{R}$ Is an Estimator with Bounded Power}
$\mathcal{R}$ in Eq.~\ref{eq:R} is a finite product over a sampled
$\mathcal{K}$, not a decision over the whole configuration space, so it is an
\emph{estimator} of robustness with quantifiable power. If a program fails a given
configuration nondeterministically with per-run probability $p$, then $N$
independent validations detect it with probability $1-(1-p)^N$; conversely,
observing $0$ failures over $N$ runs bounds the latent failure probability above by
$\approx 3/N$ at 95\% confidence (the rule of three). We use this throughout:
we report our null not as ``zero races'' but as an upper bound on the
silent-failure risk, and we characterize the estimator's power against injected
hazards as a function of input size, thread count, and repetitions
(Section~\ref{sec:results}).

\subsection{The Parallel-Correctness Gap}
Aggregating over a population of generated programs, the \emph{robust-correctness
rate} is $\overline{\mathcal{R}}=\mathbb{E}[\mathcal{R}(G)]$ and the
\emph{single-configuration rate} is $\overline{\mathcal{S}}=\mathbb{E}[\mathcal{S}(G)]$.
Their difference isolates programs that pass the developer-default check yet are
not robustly correct:
\begin{equation}\label{eq:delta}
\Delta \;=\; \Pr\!\big[\mathcal{S}(G){=}1 \ \wedge\ \mathcal{R}(G){=}0\big],
\end{equation}
the \emph{parallel-correctness gap}, with associated conditional
$\Pr[\mathcal{R}{=}0\mid\mathcal{S}{=}1]$, the \emph{silent-failure risk}: of
programs that look correct at low parallelism, the fraction actually wrong at some
higher thread count. This gap is the surface of bugs that neither a developer's
single run nor an agent's single-execution check can see.

\subsection{Verdicts and Hazard Classes}
For reporting we assign each program a mutually exclusive verdict:
\emph{robust} ($\mathcal{R}{=}1$); \emph{flaky}
($\exists k:\mathrm{corr}(G,k){=}1$ but $\mathcal{R}{=}0$, i.e.\ correct at some
thread counts or repetitions and not others---the signature of a data race);
\emph{wrong} (compiles and runs but never matches the reference); and
\emph{crash} (fails to compile or run under every configuration). Orthogonally,
we tag each program by the \emph{parallel hazard} it may exhibit, summarized in
Table~\ref{tab:hazards}. Each hazard has a characteristic differential
signature---whether the failure varies across thread counts or is invariant---that
the harness detects automatically.

\begin{table}[t]
\caption{Parallel hazard classes probed by the benchmark and their signature
under differential (thread-swept, repeated) execution. Inv.\ = the wrong result
is invariant across thread counts; \emph{varies} = the result changes with the
thread count or across repetitions.}
\label{tab:hazards}
\centering\footnotesize
\begin{tabular}{@{}p{2.5cm}p{3.1cm}l@{}}
\toprule
Class & Typical cause & signature \\
\midrule
H1 data race        & shared accumulator written without \texttt{reduction}/\texttt{atomic}  & \emph{varies} \\
H2 non-associative combine & reduction whose operator is not associative/commutative as written & \emph{varies} \\
H3 order-dependent selection & ``first''/representative element chosen without a total order & \emph{varies} \\
H4 sequential dependence & scan/recurrence parallelized despite a loop-carried dependence & invariant/varies \\
\bottomrule
\end{tabular}
\end{table}

\subsection{An Agent and Its Verification Tool}
\label{sec:agent}
The metrics above measure single-shot generation. To ask whether an \emph{agent}
can repair its own parallel code, we wrap generation in a standard
propose--observe--repair loop and vary only the \emph{tool} that produces the
observation. Given a task, the agent emits a program; a tool returns an
observation; the observation is appended to the context; the agent revises. The
loop stops when the tool reports success or a repair budget $B$ is exhausted. We
define three tool levels that isolate the value of distribution-aware
verification:
\begin{itemize}
\item \textbf{L0 (single-shot).} No tool; the first program is accepted. This is
the baseline, $\mathcal{S}$/$\mathcal{R}$ of Section~\ref{sec:metric}.
\item \textbf{L1 (execute-once).} The program is compiled and run under the
developer-default configuration $k_0$ ($\tau{=}1$); a crash or compile trace is
returned, but the program is validated at only one thread count, so any program
that runs and matches at $\tau{=}1$ is accepted. This is what an ordinary
execute-and-fix agent observes.
\item \textbf{L2 (differential verify).} The verifier of
Section~\ref{sec:metric} compiles the program and executes it across the thread
sweep with repetitions, checking each run against the serial reference $R_t$ and
returning \textsc{pass} or a specific diagnostic (a failing thread count, a
race-indicating variation across repetitions, or a crash).
\end{itemize}
Only the observation differs across levels; the model, prompt, and budget are
held fixed, so any difference in final correctness is attributable to the
\emph{tool}. We report the robust-correctness rate $\overline{\mathcal{R}}$ of
each level's \emph{final} program and the average number of repair rounds it
consumes.
